\documentclass[10pt]{article}

\usepackage[utf8]{inputenc}

\usepackage[T1]{fontenc}

\usepackage{amsmath}

\usepackage{amsfonts}

\usepackage{amssymb}

\usepackage[version=4]{mhchem}

\usepackage{stmaryrd}



\begin{document}

METРИЧЕСКИЙ OПEPATOP - см. ГИльбертово пространство.



MЕТРИЧЕСКИЙ ТЕНЗОР - дваждЫ ковариантныЙ симметричный тензор $g_{i j}(x)$, заданный в области риманова пространства с координатами $x=\left(x^{1}, x^{2}, \ldots, x^{n}\right)$, причем матрица $g_{i j}$ положительно определена: $g_{i j} T^{i} T^{j}>0$, если вектор $T \neq 0$ (принято соглашение о суммировании по повторяющимся индексам). При замене координат $x^{i} \rightarrow y^{i}(x)$ М. т. $g_{i j}$ переходит в



\[

\tilde{g}_{i j}=g_{k l}\left(\partial x^{k} / \partial y^{i}\right)\left(\partial x^{l} / \partial y^{j}\right) .

\]



М. т. иногда называют р имановой метрикой, поскольку он определяет расстояние в римановом пространстве: если задана кривая $x^{i}=x^{i}(t), a \leqslant t \leqslant b$, то ее длина



\[

s=\int_{a}^{b} d t\left[g_{i j}(x(t)) \frac{d x^{i}}{d t} \frac{d x^{j}}{d t}\right]^{1 / 2}

\]



а элемент длины $d s$ определен формулой



\[

d s=g_{i j} d x^{i} d x^{j}

\]



правая часть к-рой называется п е р в о й (основной) к в а дра т и ч о й формой. Элемент объема



\[

d V=\sqrt{g} d x^{1} \ldots d x^{n}

\]



а объем $V(U)$ области $U$ равен



\[

V(U)=\int_{U} \sqrt{g} d x^{1} \ldots d x^{n},

\]



где $g=\operatorname{det}\left\|g_{i j}\right\|$. Если существуют координаты $z^{i}$, в к-рых М. т. имеет вид $g_{i j}(z)=\delta_{i j}$, где $\delta_{i j}-$ символ Кронекера, то метрика называется е в кл и д о в о й, а сама область риманова пространства является областью евклидова пространства.



Кроме М.т. в римановом пространстве вводится еще одна независимая структура - связность, задаюшая ковариантную производную $\nabla_{k}$. М. т. называется с огласован н ы м со связностью, если он ковариантно постоянен:



\[

\nabla_{k} g_{i j}=0 .

\]



Тогда коэффициенты связности, или Кристоффеля символы, однозначно выражаются через М. т.:



\[

\Gamma_{i j}^{k}=\frac{1}{2} g^{k l}\left(\frac{\partial g_{l j}}{\partial x^{j}}+\frac{\partial g_{l i}}{\partial x^{j}}-\frac{\partial g_{i j}}{\partial x^{l}}\right) \text {. }

\]



В окрестности любой точки $x_{0}$ можно ввести нормальные (римановы) координаты такие, что $\left.\Gamma_{i j}^{k}\right|_{x=x_{0}}=0$ или $\partial g_{i j} /\left.\partial x^{k}\right|_{x=x_{0}}=0$. Тогда в этой окрестности



\[

g_{i j}=\delta_{i j}-(1 / 3) R_{i j, k l} x^{k} x^{l}+\ldots .

\]



Коэффициенты $R_{i j, k l}$ характеризуют отличие М.т. от евклидова и являются компонентами кривизны тензора. Помимо внутренних характеристик многообразия М.т. задает скалярное произведение векторов $\xi=\left(\xi^{1}, \ldots, \xi^{n}\right)$ и $\eta=\left(\eta^{1}\right.$, ..., $\left.\eta^{n}\right)$, касательных к многообразию в данной точке:



\[

(\xi, \eta)=g_{i j} \xi^{i} \eta^{j}

\]



скалярное произведение не зависит от выбора системы координат.



Понятие М. т. общеупотребительно при описании сплошной среды, при формулировке теории поля в криволинейных координатах, а особенно - в теории относительности и теории тяготения.



Лит.: [1] Л анд а у Л. Д., Л и фш и « Е. М., Теория поля, 7 изд., М., 1988; [2] Р а ш в с к и й П. К., Риманова геометрия и тензорный анализ, 3 изд., М., 1967; [3] Фок В. А., Теория пространства, времени и тяготения, 2 изд., М., 1961; [4] Дубро в и н Б. А., Нов ик о в С. П., Фо м е н ко А. Т., Современная геометрия, 2 изд., М., 1986. В. П. Павлов.



METPИЧЕСКОЕ ПРОСТРАНСТВО - см. Метрика.



358 МЕТРИЧЕСКИЙ MEХАНИКА - наука о механическом движении материальных тел и происходящих при этом взаимодействиях между ними. Под механич. движением понимают изменение с течением времени взаимного положения тел или их частиц в пространстве. В природе - это движение небесных тел, колебания земной коры, воздушные и морские течения и т. П., а в технике - движения различных летательных аппаратов и транспортных средств, частей двигателей, машин и механизмов, деформации элементов различных конструкций и сооружений, движения жидкостей и газов и многое другое. Рассматриваемые в М. взаимодействия представляют собой те действия тел друг на друга, результатами к-рых являются изменения скоростей точек этих тел или их деформации, напр. притяжения тел по закону всемирного тяготения, взаимные давления соприкасающихся тел, воздействия частиц жидкости или газа друг на друга и на движущиеся в них тела.



Под М. обычно понимают так наз. к л а с с и ч е с к у ю м е х а н и ку, в основе к-рой лежат Ньютона законы механи$\kappa и$, а предметом ее изучения являются движения любых материальных тел (кроме элементарных частиц), совершаемые со скоростями, малыми по сравнению со скоростью света. Движение тел со скоростями порядка скорости света рассматриваются в относительности теории, а внутриатомные явления и движение элементарных частиц изучаются в квантовой механике.



При изучении движения материальных тел в М. вводят ряд абстрактных понятий, отражающих те или иные свойства реальных тел; ими являются:



\begin{enumerate}

  \item материальная точка - объект пренебрежимо малых размеров, имеющий массу; это понятие применимо, когда тело движется поступательно или когда в изучаемом движении можно пренебречь врашением тела вокруг его центра масс.



  \item Абсолютно твердое тело-тело, расстояние между двумя любыми точками к-рого всегда остается неизменным; это понятие применимо, когда можно пренебречь деформацией тела.



  \item Сплошная изменяемая среда; это понятие применимо, когда при изучении движения изменяемой среды (деформируемого твердого тела, жидкости, газа) можно пренебречь молекулярной структурой среды. При изучении сплошных сред прибегают к следующим абстракциям, отражающим при данных условиях наиболее существенные свойства соответствующих реальных тел: идеально упругое тело, пластич. тело, идеальная жидкость, вязкая жидкость, идеальный газ и др.



\end{enumerate}



В соответствии с этим М. разделяют на М. материальной точки, М. системы материальных точек, М. абсолютно твердого тела и М. сплошной среды; последняя в свою очередь подразделяется на теорию упругости, теорию пластичности, гидродинамику, аэродинамику, газовую динамику и др. В каждом из этих подразделов в соответствии с характером решаемых задач выделяют: статику - учение о равновесии тел под действием сил, кинематику - учение о геометрич. свойствах движения тел и динамику - учение о движении тел под действием сил. Изучение основных законов и принципов, к-рым подчиняется механич. движение тел, и вытекающих из этих законов и принципов общих теорем и уравнений составляет содержание так наз. общ е й, или т е о р е т и ч е с к о й, м е ха н и к и. Разделами М., имеющими самостоятельное значение, являются также колебаний и волн теория, теория устойчивости равновесия и устойчивости движения, теория гироскопа, математическая теория, механика тел переменной массы, теория автоматич. регулирования, теория удара и др. Важное место в М., особенно в М. сплошных сред, занимают экспериментальные исследования, проводимые с помощью разнообразных механич., оптич., электрич. и других физич. методов и приборов. М. тесно связана со многими другими разделами физики. Ряд понятий и методов М. при соответствующих обобщениях находит приложение в оптике, статистич. физике, квантовой М., электродинамике, теории относительности и др. (см., напр., Действие, Лагранжа функция, Лагранжа уравнения в общей механике). Кроме того, при решении





\end{document}